%
% uebersicht.tex -- Uebersicht ueber die Seminar-Arbeiten
%
% (c) 2022 Prof Dr Andreas Mueller, OST Ostschweizer Fachhochschule
%
\chapter*{Übersicht}
\fancyhead[RE]{}
\fancyhead[LO]{Übersicht}
\label{buch:uebersicht}
Im zweiten Teil kommen die Teilnehmer des Seminars selbst zu Wort.
Die im ersten Teil dargelegten mathematischen Methoden und
grundlegenden Modelle werden dabei verfeinert, verallgemeinert
und auch numerisch überprüft.

%
% Basel-Problem
%
Das Basel-Problem hat die besten Mathematiker des 17.~und 18.~Jahrhunderts
herausgefordert.
\index{Basel-Problem}%
Erst Leonhard Euler fand die korrekte Summe der reziproken Quadrate der
\index{Euler, Leonhard}%
natürlichen Zahlen.
\emph{Jan Gachnang}
und
\emph{Etienne Schafflützel}
präsentieren in Kapitel~\ref{chapter:basel} Eulers klassische Lösung,
die allerdings modernen Standards nicht genügt.
Strenge Beweise werden ebenfalls dargelegt.

%
% Konstruktion
%
Mit komplexen Zahlen lassen sich geometrische Konstruktionen in
\index{Konstruktion, geometrisch}%
der Ebene besonders einfach beschreiben.
\emph{Timon Gnehm}
\index{Gnehm, Timon}%
\index{Timon Gnehm}%
und
\emph{Pascal Hirzel}
\index{Pascal Hirzel}%
\index{Hirzel, Pascal}%
stellen in Kapitel~\ref{chapter:konstruktion} den Satz von Napoleon
\index{Satz!von Napoleon}%
vor.

%
% Produktdarstellungen
%
Polynome sind bis auf einen Vorfaktor durch die Nullstellen
festgelegt.
Ob dies auch für beliebige ganze Funktionen gilt, diskutieren
\emph{Simon Köpfli}
\index{Kopfli, Simon@Köpfli, Simon}%
\index{Simon Kopfli@Simon Köpfli}%
und
\emph{Florian von Wyl}
\index{von Wyl, Florian}%
\index{Florian von Wyl}%
in Kapitel~\ref{chapter:produkt}.
Da es ganze Funktionen mit unendlich vielen Nullstellen gibt, muss
man dazu unendlichen Produkten einen Sinn geben.
\index{Produkt, unendlich}%
Dabei treten zunächst Konvergenzfragen als Hindernis auf, die durch
die weierstraßschen Elementarfaktoren beantwortet werden.
Dies genügt aber nicht.
Die Exponentialfunktion $e^{g(z)}$ einer ganzen Funktion $g(z)$ hat
keine Nullstellen, die möglicherweise verschiedenen Funktionen $f(z)$
und $e^{g(z)}f(z)$ haben daher die gleichen Nullstellen, es können
aber im Allgemeinen nicht beide Funktionen Polynome sein.

%
% Julia-Mengen und Mandelbrot-Menge
%
Die Theorie der dynamischen System befasst sich mit dem Verhalten
von Iterationsfolgen.
Holomorphe Funktionen sind in diesem Zusammenhang besonders interessant,
da sie zu winkeltreuen Abbildungen der komplexen Ebene gehören.
\emph{Nicola Dall'Acqua} zeigt in Kapitel~\ref{chapter:julia}, wie
\index{Nicola Dall'Acqua}%
\index{Dall'Acqua, Nicola}%
die Julia- und Fatou-Mengen Punkte der komplexen Ebene anhand des
\index{Julia-Menge}%
\index{Fatou-Menge}%
Konvergenzverhaltens der zugehörigen Iterationsfolge klassifizieren.
Für holomorphe Funktionen entstehen filigrane selbstähnlichen Strukturen,
die diesen Mengen auch eine besondere ästhetische Qualität geben.
Die Mandelbrot-Menge im Parameterraum für quadratische Abbildungen
\index{Mandelbrot-Menge}%
wird ebenfalls eingeführt.

%
% Möbius-Transformationen
%
Die Bewegungen der ebenen euklidischen Geometrie können in der 
komplexen Ebene breiter gefasst werden.
Indem die komplexe Ebene zur riemannschen Zahlenkugel erweitert wird,
\index{Zahlenkugel, riemannsch}%
können gebrochen rationale Transformationen, auch Möbius-Transformationen
\index{Mobius-Transformation@Möbius-Transformation}%
genannt, als Transformationsgruppe verwendet werden.
Sie sind zwar nicht mehr längentreu, aber immer noch winkeltreu.
\emph{Tom Dawson}
\index{Tom Dawson}%
\index{Dawson, Tom}%
zeigt in Kapitel~\ref{chapter:moebius}, wie man man Möbius-Transformationen
finden, analysieren und damit rechnen kann.
Insbesondere zeigt sich, dass die Zusammensetzung von Möbius-Transformationen
durch die Matrixmultiplikation beschrieben werden kann.

%
% Complex Step Derivative
%
In der reellen numerischen Analysis ist die Bestimmung zuverlässiger
Werte für die Ableitung ein besonderes Problem.
Ursache dafür sind die unterschiedlichen Grössenordnungen von Koordinaten
und Koordinatendifferenzen.
\index{finite Differenz}%
\index{Differenz, finite}%
Bei der Berechnung der Ableitung mit finiten Differenzen wird die
Genauigkeit durch starke Auslöschung reduziert.
\index{Ausloschung@Auslöschung}%
\emph{Selina Malacarne} zeigt in Kapitel~\ref{chapter:step}, wie die
\index{Selina Malacarne}%
\index{Malacarne, Selina}%
Ableitung in imaginärer Richtung für holomorphe Funktionen auf natürliche
Art die Trennung von Koordinaten und Koordinatendifferenz ermöglicht
und so das Problem der Auslöschung umgeht.
Eingeführt wird die Idee mithilfe der dualen Zahlen, die als eine
\index{duale Zahlen}%
Vorstufe der Methoden der Nichtstandard-Analysis zur Berechnung der
\index{Nichtstandard-Analysis}%
Ableitung angesehen werden können.


